% !TEX root = ../../../main.tex
% 所属：第四部分：Fourier 分析 / 第4章：$\R$上的Fourier变换
% 节：4.5 Heisenberg 不确定性原理
% 正文按原文件顺序保存；由 ../chapter4.tex 调用。
\section{Heisenberg 不确定性原理}

\begin{theorem}\label{thm:heisenberg-uncertainty}
	假设 $\psi\in\mathcal S(\R)$，并且满足标准化条件
	\[
	\int_{-\infty}^{\infty}\abs{\psi(x)}^2\dd x=1.
	\]
	则
	\[
	\left(\int_{-\infty}^{\infty}x^2\abs{\psi(x)}^2\dd x\right)
	\left(\int_{-\infty}^{\infty}\xi^2\abs{\widehat\psi(\xi)}^2\dd\xi\right)
	\geq \frac{1}{16\pi^2},
	\]
	且等号成立当且仅当
	\[
	\psi(x)=Ae^{-Bx^2},
	\]
	其中 $B>0$，且
	\[
	\abs{A}^2=\sqrt{\frac{2B}{\pi}}.
	\]
	事实上，对任意 $x_0,\xi_0\in\R$，都有
	\[
	\left(\int_{-\infty}^{\infty}(x-x_0)^2\abs{\psi(x)}^2\dd x\right)
	\left(\int_{-\infty}^{\infty}(\xi-\xi_0)^2
	\abs{\widehat\psi(\xi)}^2\dd\xi\right)
	\geq \frac{1}{16\pi^2}.
	\]
\end{theorem}

\begin{proof}
	将 $\psi(x)$ 替换为
	\[
	e^{-2\pi i\xi_0x}\psi(x+x_0),
	\]
	再通过变量替换，便可由第一个不等式推出第二个不等式。因此只需证明 $x_0=\xi_0=0$ 的情形。
	
	由标准化条件以及 $\psi$ 的迅速下降性，分部积分可得
	\[
	\begin{aligned}
		1
		&=\int_{-\infty}^{\infty}\abs{\psi(x)}^2\dd x\\
		&=-\int_{-\infty}^{\infty}x\frac{\dd}{\dd x}
		\abs{\psi(x)}^2\dd x\\
		&=-\int_{-\infty}^{\infty}
		\left(x\psi'(x)\overline{\psi(x)}
		+x\overline{\psi'(x)}\psi(x)\right)\dd x.
	\end{aligned}
	\]
	最后一个等式成立是因为 $\abs{\psi}^2=\overline\psi\psi$。因而由 Cauchy--Schwarz 不等式可得
	\[
	\begin{aligned}
		1
		&\leq 2\int_{-\infty}^{\infty}
		\abs{x}\abs{\psi(x)}\abs{\psi'(x)}\dd x\\
		&\leq 2
		\left(\int_{-\infty}^{\infty}x^2\abs{\psi(x)}^2\dd x\right)^{1/2}
		\left(\int_{-\infty}^{\infty}\abs{\psi'(x)}^2\dd x\right)^{1/2}.
	\end{aligned}
	\]
	由 Fourier 变换的性质以及 Plancherel 公式可知
	\[
	\int_{-\infty}^{\infty}\abs{\psi'(x)}^2\dd x
	=4\pi^2\int_{-\infty}^{\infty}
	\xi^2\abs{\widehat\psi(\xi)}^2\dd\xi.
	\]
	由此立即得到定理中的不等式。
	
	如果等号成立，则在应用 Cauchy--Schwarz 不等式时等号也成立，并且相应内积为负实数。因此存在实常数 $\beta<0$，使得
	\[
	\psi'(x)=\beta x\psi(x).
	\]
	该方程的解为
	\[
	\psi(x)=Ae^{\beta x^2/2}.
	\]
	令 $\beta=-2B$，则 $B>0$，从而
	\[
	\psi(x)=Ae^{-Bx^2}.
	\]
	再由标准化条件，
	\[
	1=\abs{A}^2\int_{-\infty}^{\infty}e^{-2Bx^2}\dd x
	=\abs{A}^2\sqrt{\frac{\pi}{2B}},
	\]
	故
	\[
	\abs{A}^2=\sqrt{\frac{2B}{\pi}}.
	\]
	反之，这类函数代入上述推导时处处取等号。
\end{proof}
